\documentclass[11pt,compress,t,notes=noshow, xcolor=table]{beamer}

\input{../../style/preamble}
\input{../../latex-math/basic-math}
\input{../../latex-math/basic-ml}

\newcommand{\zin}[2]{z_{\text{in},{#1}}^{(#2)}}
\newcommand{\zout}[2]{z_{\text{out},{#1}}^{(#2)}}

\title{Optimization in Machine Learning}

\begin{document}

\comment[NOTE]{BB}{wir sollten dringend das hier mal durchsehen und uns daran orientieren https://ruder.io/optimizing-gradient-descent/}
\comment[NOTE]{BB}{wir sollten generell dieses gesamt subcha machen zu normalied GD --> https://jermwatt.github.io/machine_learning_refined/notes/3_First_order_methods/3_9_Normalized.html}
\comment[NOTE]{BB}{vielleicht sollten wir mal definieren was first order überhaupt heisst. und welche first order verfahren es gibt -- egal ob wir die behandeln}

\titlemeta{
First order methods
}{
Gradient descent
}{
figure_man/linesearch.png
}{
\item Iterative Descent / Line Search
\item Descent directions
\item Gradient descent
\item Pseudo-residuals
\item ERM with GD
}

\begin{framei}{Iterative Descent}
\item \textbf{Goal}: For $f: \R^d \to \R$ continuously differentiable, solve
$$\argmin_{\xv \in \R^d} f(\xv)$$
\item Usually no closed-form solution; construct iterates $\xv^{[0]}, \xv^{[1]}, \dots$ using only \textbf{local} information $f(\xv)$, $\nabla f(\xv)$
\item Metaphor: $f$ = height of a mountain at coordinates $(x_1, x_2)$; in dense fog, we cannot see the valley, only sense the local slope
\item \textbf{Central idea}: Repeat
\splitVCC[0.6]{
\begin{enumerate}
\item At current $\xv \in \R^d$, find a \textbf{descent direction} $\mathbf{d} \in \R^d$
\item Update $\xv \leftarrow \xv + \alpha \mathbf{d}$, with step size $\alpha$ s.t. $f$ sufficiently decreases (\textbf{step size control})
\end{enumerate}}{
\imageC{figure_man/linesearch.png}
\begin{footnotesize} ``Walking down the hill through local steps, towards the valley''\end{footnotesize}}
\end{framei}


\begin{framei}{Iterative Descent}
\item \textbf{Definition:} $\mathbf{d}\in \R^d$ is a \textbf{descent direction} in $\xv$ if
$$D_\mathbf{d} f(\xv) = \nabla f(\xv) ^T \mathbf{d} < 0 \text{ (neg. directional derivative)}$$
\begin{figure}
\imageC[0.4]{figure/descent_direction.pdf}
\begin{footnotesize}
Angle between $\nabla \fx$ and $\mathbf{d}$ must be $> 90^{\circ}$.
\end{footnotesize}
\end{figure}
\end{framei}


\begin{framev}{Iterative Descent}
\begin{algorithm}[H]
\caption{Iterative Descent / Line search}
\begin{algorithmic}[1]
\State Starting point $\xv^{[0]} \in \R^d$
\While {Stopping criterion not met}
\State Calculate a descent direction $\mathbf{d}^{[t]}$ for current $\xv^{[t]}$
\State Find $\alpha^{[t]}$ s.t. $f(\xv^{[t]} + \alpha^{[t]} \mathbf{d}^{[t]}) < f(\xv^{[t]})$
\State Update $\xv^{[t + 1]} = \xv^{[t]} + \alpha^{[t]} \mathbf{d}^{[t]}$
\EndWhile
\end{algorithmic}
\end{algorithm}
Key questions:
\begin{itemizeM}
\item How to choose $\mathbf{d}^{[t]}$ (now)
\item How to choose $\alpha^{[t]}$ (later)
\end{itemizeM}
\vfill
\begin{footnotesize}
NB: Above procedure also called \textbf{line search}, but term is ambiguous: can serve as umbrella term for the whole scheme or refer only to Step 4.
\end{footnotesize}
\end{framev}


\begin{framei}{Gradient Descent}
\item Locally, $\nabla \fx$ is the direction of steepest ascent, $-\nabla \fx$ of steepest descent
\item \textbf{Gradient descent (GD)}: iterative descent with $\mathbf{d} = -\nabla \fx$:
$$\xv^{[t+1]} = \xv^{[t]} - \alpha^{[t]} \nabla f(\xv^{[t]})$$
\item \textbf{Gradient ascent (GA)}: maximize $f$ analogously via:
$$\xv^{[t+1]} = \xv^{[t]} + \alpha^{[t]} \nabla f(\xv^{[t]})$$
\splitV{\imageC[0.8]{figure_man/gradient-descent-example_1.png}}{\imageC[0.9]{figure_man/gradient-descent-example_2.png}}
\begin{center}\begin{footnotesize}
GD for $f(x_1, x_2) = - \sin(x_1) \cdot \varphi(x_2)$, with $\varphi$ the density of $\normal(\pi/2, 0.8^2)$.
\end{footnotesize}\end{center}
\end{framei}


\comment[NOTE]{BB}{bit unclear whether we should keep the GD-ERM stuff here or do it later....?
for now maybe fine} 
\begin{framei}{GD for empirical risk minimization}
\item Let $\D = \Dset$ and consider ERM:
$$\argmin_{\thetav} \risket = \argmin_{\thetav} \sumin \Lxyit$$
\item \textbf{Gradient} of the empirical risk:
$$\nabla_{\thetav} \risket = \sumin \pd{\Lxyit}{\thetav}$$
\item \textbf{GD update rule}:
$$\thetav^{[t+1]} = \thetav^{[t]} - \alpha^{[t]} \nabla_{\thetav} \risket\big|_{\thetav^{[t]}}$$
\end{framei}


\begin{framei}[fs=footnotesize]{GD for ERM: Pseudo-Residuals}
\item By the chain rule, the (neg.) gradient of the ERM problem can be expressed as:
$$
- \nabla_{\thetav} \risket = \sumin \underbrace{\left( - \pd{\Lxyi}{f} \right)}_{\text{pseudo-residual } \tilde r^{(i)}(f)} \pd{\fxit}{\thetav}
$$
\item \textbf{Pseudo-residuals} tell us how to distort $\fxi$ to achieve greatest decrease of $\Lxyi$ (best pointwise update)
\item $\pd{\fxit}{\thetav}$ tells us how to modify $\thetav$ accordingly and wiggle model output
\item GD step sums up these modifications across all observations $i$
\splitVCC{
\textbf{NB:} Pseudo-residuals
$\tilde{r}\left( f \right)$
match usual residuals for L2 loss:
\begin{align*}
\tilde{r}(f) = - \pd{\Lxy}{f} & = - \frac{1}{2} \left. \pd{(y - f)^2}{f} \right|_{f = \fx}\\
& = y - \fx
\end{align*}
}
{\imageC{figure_man/pseudo_residual_1.png}}
\end{framei}


\begin{framei}{Optimize LS linear regression with GD}
\item Let $\D = \Dset$ and minimize
$$\risket = \sumin (\thetav^\top \xi - \yi)^2$$
{\footnotesize
\textbf{NB:} For illustration, we use GD even though closed-form solution exists. GD-like (more adv.) approaches like this MIGHT make sense for large data, though.}
\item\textbf{Gradient: } $$\nabla_{\thetav}\risket = \frac{\partial \risket}{\partial \thetav} = -2 \sumin (\yi - \thetav^{\top} \xi) \xi$$
\imageC[0.8]{figure/gradient_descent_lm.pdf}
\end{framei}


\begin{framei}{GD for logistic regression}
\item Let $\D = \Dset$ with $\yi \in \{0, 1\}$ and minimize the cross-entropy risk
$$\risket = - \sumin \left[\yi \log \pi^{(i)} + (1 - \yi) \log(1 - \pi^{(i)})\right]$$
where $\pi^{(i)} := \pixit = s(\thetav^\top \xi)$, \quad $s(z) = \tfrac{1}{1 + \exp(-z)}$
\item \textbf{Gradient:}
$$\nabla_{\thetav} \risket = \sumin (\pi^{(i)} - \yi) \xi$$
\imageC[0.8]{figure/gradient_descent_logistic.pdf}
\end{framei}


\begin{framei}{ERM for NN with GD}
\item Consider $y = x_1^2 + x_2^2$ and $\D = \Dset$
\vspace{-0.75cm}
\imageC[0.5]{figure/gradient_descent_NN_0.pdf}
\item Fit a neural network $\fxt$ with two hidden layers (16 ReLU units each) by minimizing the empirical risk
$$\risket = \sumin \left(\fxt - \yi\right)^2$$

\end{framei}


\begin{framei}[fs=footnotesize]{Backpropagation}
\item GD requires $\nabla_{\thetav} \risket$; for NNs computed via \textbf{backpropagation} (BP)
\item Weight $\theta$ affects $L$ through \textbf{all} paths $\theta \to z_{\text{in}} \to z_{\text{out}} \to \dots \to f \to L$ in the network; chain rule: $\pd{L}{\theta}$ = sum over these paths of products of \textbf{local} derivatives
\item BP: \textbf{forward pass} computes all unit values; \textbf{backward pass} accumulates the path sums from output to input, reusing shared factors
\item \textbf{Example} (2 hidden layers with 2 units each, activation $\sigma$):
let $\ L = \tfrac{1}{2}(y - f(x))^2$ and
$\ \zin{1}{1} = \theta_{11}^{(1)} x_1 + \theta_{12}^{(1)} x_2$,
$\ \zout{1}{1} = \sigma(\zin{1}{1})$, etc.
$$
\pd{L}{\theta_{11}^{(3)}} = \underbrace{\pd{L}{f}}_{f-y} \cdot \underbrace{\pd{f}{\theta_{11}^{(3)}}}_{{\zout{1}{2}}}, \qquad
\pd{L}{\theta_{11}^{(2)}} = \underbrace{\pd{L}{f}}_{f-y} \cdot \underbrace{\pd{f}{\zout{1}{2}}}_{\theta_{11}^{(3)}} \cdot \underbrace{\pd{\zout{1}{2}}{\zin{1}{2}}}_{\sigma'(\zin{1}{2})} \cdot \underbrace{\pd{\zin{1}{2}}{\theta_{11}^{(2)}}}_{\zout{1}{1}}, \qquad
\dots
$$
\vspace{-1cm}
\imageC[0.45]{figure/nn_architecture.pdf}
{\scriptsize
Note: Shown weights have a unique path to $f$ (single products); layer-1 weights reach $f$ via both layer-2 units $\Rightarrow$ sum of two products. Biases omitted; actual NN: 16 units/layer.
}
\end{framei}


\begin{framei}[fs=footnotesize]{Excursus: BP error signals}
\item Vector notation for hidden layers $l = 1, \dots, m$, with weight matrices $\bm{\Theta}^{(l)} = (\theta_{jk}^{(l)})_{j,k}$ and $\sigma$ applied elementwise:
$$\bm{z}_{\text{in}}^{(l)} = \bm{\Theta}^{(l)} \bm{z}_{\text{out}}^{(l-1)}, \; \text{ i.e., } \zin{j}{l} = \sum_k \theta_{jk}^{(l)} \zout{k}{l-1}, \qquad
\bm{z}_{\text{out}}^{(l)} = \sigma(\bm{z}_{\text{in}}^{(l)})$$
\item Input, linear output, and L2 loss:
$$\bm{z}_{\text{out}}^{(0)} = \xv, \qquad f = \bm{\theta}^{(m+1) \top} \bm{z}_{\text{out}}^{(m)}, \qquad L = \tfrac{1}{2} (y - f)^2$$
\item \textbf{Error signal} of layer $l$ = sensitivity of $L$ w.r.t. unit inputs:
$$\delta_j^{(l)} := \pd{L}{\zin{j}{l}}, \qquad \text{at the output: } \delta^{(m+1)} = \pd{L}{f} = f - y$$
\item $\theta_{jk}^{(l)}$ enters $L$ only via $\zin{j}{l}$, and $\pd{\zin{j}{l}}{\theta_{jk}^{(l)}} = \zout{k}{l-1}$ (see def.\ above), so \textbf{all} gradients follow directly from the error signals:
$$\pd{L}{\theta_{jk}^{(l)}} = \pd{L}{\zin{j}{l}} \cdot \pd{\zin{j}{l}}{\theta_{jk}^{(l)}} = \delta_j^{(l)} \zout{k}{l-1}, \qquad \text{in matrix form } \pd{L}{\bm{\Theta}^{(l)}} = \bm{\delta}^{(l)} \big( \bm{z}_{\text{out}}^{(l-1)} \big)^\top$$
\item Missing piece: computing $\bm{\delta}^{(l)}$ for all layers $\to$ next slide
\end{framei}


\begin{framei}[fs=footnotesize]{Excursus: BP error signal recursion}
\item $\zin{j}{l}$ enters $L$ only via the \textbf{next} layer's inputs, which depend on it as
$$\zin{k}{l+1} = \sum_{j'} \theta_{kj'}^{(l+1)} \zout{j'}{l} = \sum_{j'} \theta_{kj'}^{(l+1)} \sigma(\zin{j'}{l})
\quad \Rightarrow \quad \pd{\zin{k}{l+1}}{\zin{j}{l}} = \theta_{kj}^{(l+1)} \sigma'(\zin{j}{l})$$
\item Chain rule, summing over all units $k$ of layer $l+1$ (= sum over paths), yields the \textbf{backward recursion} for $l = m, \dots, 1$:
\begin{align*}
\delta_j^{(l)} &= \sum_k \pd{L}{\zin{k}{l+1}} \cdot \pd{\zin{k}{l+1}}{\zin{j}{l}} = \sigma'(\zin{j}{l}) \sum_k \theta_{kj}^{(l+1)} \delta_k^{(l+1)}\\
\text{i.e., } \quad \bm{\delta}^{(l)} &= \big( \bm{\Theta}^{(l+1) \top} \bm{\delta}^{(l+1)} \big) \odot \sigma'(\bm{z}_{\text{in}}^{(l)}) \qquad (\odot = \text{elementwise product})
\end{align*}
started at $\delta^{(m+1)} = f - y$, reading $\bm{\Theta}^{(m+1)} = \bm{\theta}^{(m+1) \top}$ for the output layer
\item Together: \textbf{forward pass} computes all $\bm{z}_{\text{in}}^{(l)}, \bm{z}_{\text{out}}^{(l)}$; \textbf{backward pass} computes all $\bm{\delta}^{(l)}$, then all gradients $\bm{\delta}^{(l)} \big( \bm{z}_{\text{out}}^{(l-1)} \big)^\top$
\item Cost: one forward + one backward pass $\Rightarrow$ full gradient costs a small constant multiple ($\approx 2$--$4\times$) of \textbf{one} evaluation of $f$; this factor does not grow with the number of weights $d$ (compare finite differences: $d + 1$ evaluations)
\item To see the factor: forward, backward, and readout ($\bm{\Theta}^{(l)} \bm{z}_{\text{out}}^{(l-1)}$, $\bm{\Theta}^{(l+1) \top} \bm{\delta}^{(l+1)}$, $\bm{\delta}^{(l)} \big( \bm{z}_{\text{out}}^{(l-1)} \big)^\top$) each touch every weight once $\Rightarrow$ $\approx 3d$ vs.\ $d$ multiply-adds for $f$
\end{framei}


\begin{framev}[align=top]{ERM for NN with GD}
\only<1>{
After $10$ iters of GD:
\splitV{
\imageFixed{0cm}{1.5cm}[1]{figure/gradient_descent_NN_10_surface.pdf}}
{\imageFixed{6cm}{2.5cm}[1]{figure/gradient_descent_NN_10_history.pdf}}
}
\only<2>{
After $100$ iters of GD (note the zig-zag pattern of the risk after iter. 30):
\splitV{
\imageFixed{0cm}{1.5cm}[1]{figure/gradient_descent_NN_100_surface.pdf}}
{\imageFixed{6cm}{2.5cm}[1]{figure/gradient_descent_NN_100_history.pdf}}
}
\only<3>{
After $300$ iters of GD (note the zig-zag pattern of the risk after iter. 30):
\splitV{
\imageFixed{0cm}{1.5cm}[1]{figure/gradient_descent_NN_300_surface.pdf}}
{\imageFixed{6cm}{2.5cm}[1]{figure/gradient_descent_NN_300_history.pdf}}
}
\end{framev}

\endlecture
\end{document}
